\documentclass[journal]{IEEEtran}
\usepackage{amsmath,amssymb,booktabs,graphicx,url,microtype}
\usepackage{tikz}
\usetikzlibrary{arrows.meta,positioning}
\usepackage[hidelinks]{hyperref}

\title{When Should a Vehicle Move to Learn the Channel? Decision-Triggered Active Self-Calibration for PC-FMCW-Informed Vehicular Optical Planning}
\author{Panagiota Grosdouli}

\begin{document}
\maketitle

\begin{abstract}
Communication-aware vehicle planners may benefit from adapting uncertain link models online, but aggressive information seeking can distort motion when additional knowledge would not change the selected trajectory. We study a decision-triggered active self-calibration strategy for a modeled directional optical link embedded in receding-horizon vehicular planning. Five planners share identical candidate generation, target prediction, and hard safety filtering: a fixed nominal model, passive Bayesian calibration, unconditional information seeking, decision-triggered information seeking, and a nondeployable true-parameter reference. A predeclared 27-setting development sweep on 20 seeds selected one setting before 50 untouched confirmatory seeds were opened. Across six declared scenarios, 1,500 confirmatory episodes and 45,000 timesteps recorded zero collisions, zero no-candidate steps, and zero static-clearance violations under the sampled benchmark. Unconditional information seeking increased cumulative decision regret by 0.083290 relative to passive calibration, whereas decision triggering recovered 0.082211 of that penalty and reduced probe fraction by 0.037444; both effects survived Holm correction. However, decision-triggered calibration did not establish superiority over passive calibration, and in the decision-critical scenario it probed without lowering regret. A distance-only development ablation eliminated nontrivial probing, while a measured vehicular visible-light study showed essentially null directional predictive gain. The evidence therefore supports decision relevance as a mechanism for suppressing unnecessary probing, not as proof that active calibration outperforms passive Bayesian updating or that the modeled optical parameters are physically calibrated.
\end{abstract}

\begin{IEEEkeywords}
active calibration, communication-aware motion planning, dual control, vehicular optical communication
\end{IEEEkeywords}

\section{Introduction}
Connected and automated vehicles increasingly rely on communication quality as part of the decision context rather than as a passive networking metric. Communication-aware motion planning is well established: robot trajectories have been optimized to avoid poor connectivity regions \cite{ghaffarkhah2011}, modern radio-map estimates have been integrated directly into robot navigation \cite{gordon2026}, and vehicular trajectory planning has explicitly targeted communication quality-of-service requirements \cite{ullah2025}. Integrated sensing and communication (ISAC) further couples sensing, communication, and transportation objectives \cite{cheng2022,jin2026}. These developments motivate a harder question that appears when the communication model itself is uncertain: should vehicle motion be chosen partly to \emph{learn} the link?

That question belongs to the broader dual-control problem. A control action can simultaneously direct the physical system and change what will be learned about uncertain dynamics or parameters \cite{barshalom1974,mesbah2018}. Active learning can therefore improve future decisions, but probing has an immediate cost. In constrained vehicular planning, that cost is especially important because an informative motion can still be strategically poor even when it is physically safe.

The central distinction in this paper is between \emph{parameter uncertainty} and \emph{decision-relevant uncertainty}. A posterior over channel parameters may remain broad while every plausible hypothesis selects the same safe trajectory. In that case, reducing posterior entropy can have little decision value. Conversely, a smaller uncertainty region may matter greatly if plausible hypotheses disagree about which candidate trajectory is best. This motivates the study question: \emph{When should safe ego motion also be used as an experiment for learning uncertain modeled optical-link parameters because the information can change the downstream motion decision?}

We study this question in a controlled PC-FMCW-informed analytical simulation. The latent optical quantities are simulator parameters, not measured calibration constants. Five planners, C0--C4, share the same candidate lattice, causal target prediction, and hard feasibility filters. They differ only in how they represent and exploit uncertainty in the modeled link. The experimental protocol separates development and confirmation, fixes the statistical unit at the independent seed, and retains null or negative findings rather than retuning them away.

The paper makes four bounded contributions:
\begin{enumerate}
\item It implements an explicit \emph{decision-relevance gate} that enables information seeking only when posterior hypotheses disagree about the belief-optimal trajectory and that disagreement carries nonzero expected decision regret.
\item It provides a causal C0--C4 decomposition that separates fixed nominal planning, passive Bayesian calibration, unconditional active information seeking, decision-triggered information seeking, and a model-relative true-parameter reference while preserving one common hard-safety interface.
\item It evaluates the mechanism under a frozen seed-level confirmatory protocol with paired bootstrap intervals, paired Wilcoxon tests, effect sizes, and within-family Holm correction. The primary result is mixed: C3 strongly improves over unconditional C2, but does not establish superiority over passive C1.
\item It bounds interpretation with a distance-only mechanism ablation and separately scoped measured-data studies. The measured vehicular visible-light comparison is essentially null for directional predictive gain, and the 5G V2X study is not relabeled as optical calibration.
\end{enumerate}

No universal priority claim is made for dual control, informative motion, active calibration, communication-aware planning, ISAC-to-planning coupling, or optical communication-aware control. The novelty question addressed here is narrower: whether \emph{downstream decision disagreement} is a useful trigger for active calibration in this particular safe vehicular planning architecture.

\section{Related Work}
\subsection{Dual Control and Active Uncertainty Reduction}
The classical dual effect describes controls that influence both physical state and future uncertainty \cite{barshalom1974}. Modern reviews of stochastic model predictive control (MPC) emphasize the same exploration--exploitation coupling and the computational difficulty of exact dual control \cite{mesbah2018}. Approximate Bayesian dual-control formulations similarly reason about how actions change future information \cite{klenske2016}. Online experiment-design MPC has explicitly introduced excitation during normal control operation and has also emphasized avoiding unnecessary excitation once adequate information has been acquired \cite{heirung2015}.

Recent work has brought these ideas closer to robotics. Li \emph{et al.} formulate dual control of exploration and exploitation with active learning for auto-optimization \cite{li2025}. Hu \emph{et al.} develop shielding-aware dual control for safe interactive motion planning and demonstrate it in autonomous-vehicle experiments \cite{hu2024}. Pashupathy \emph{et al.} combine active estimation and path planning for robotic assembly \cite{pashupathy2026}. These studies establish both the relevance of dual control to robotic motion and the importance of safety-aware exploration.

The present method is not an exact dual controller. It uses a finite candidate lattice, a discrete Bayesian grid, a short-horizon predictive-information proxy, and a receding-horizon gate. Its specific design choice is to condition active information seeking on \emph{decision disagreement and expected decision regret}, rather than on posterior spread or information gain alone.

\subsection{Informative Motion and Calibration}
Informative path planning optimizes robot trajectories to collect measurements that improve a model of the environment \cite{hitz2017,popovic2020}. Calibration research likewise exploits motion geometry: observability-aware LiDAR--IMU calibration explicitly identifies informative and non-identifiable directions \cite{lv2022}, while optimal-experiment-design trajectory planning has been used to generate calibration motions for robotic camera systems \cite{amersdorfer2025}. These works show that robot motion can itself be an experiment.

Our distinction is downstream. The objective is not to maximize calibration quality per se. The primary endpoint is trajectory-selection regret under the true simulator parameter, and the active mechanism is penalized when it learns parameters without improving the decision.

\subsection{Communication-Aware and ISAC-Aware Planning}
Communication-aware planning has a long history in mobile networks \cite{ghaffarkhah2011}. Recent work uses online radio-map estimation to build service-aware risk maps for robot motion \cite{gordon2026}, while autonomous-vehicle trajectory planning has been designed around target communication QoS \cite{ullah2025}. In vehicular ISAC, sensing and communication are increasingly treated jointly \cite{cheng2022}; planning-oriented ISAC explicitly connects sensing uncertainty and resource allocation to autonomous-vehicle motion planning \cite{jin2026}. Optical communication constraints have also been embedded directly in nonlinear MPC for aerial robots \cite{silano2025}.

Accordingly, this paper does not claim that communication can influence motion, that ISAC can interact with planning, or that optical-link quality can enter a control objective. The narrower question is when uncertainty in a modeled directional link should justify \emph{intentional probing motion}.

\subsection{Vehicular Optical Links and Channel Modeling}
Optical V2V communication has been demonstrated with LED transmitters and camera receivers under outdoor and driving conditions \cite{takai2014}. Vehicular visible-light communication (V-VLC) models account for distance, lateral displacement, transmitter geometry, weather, and receiver properties \cite{karbalayghareh2020}. Measurement-driven work has used machine learning for V-VLC path-loss modeling \cite{turan2021}, characterized reflected non-line-of-sight channels \cite{turan2022nlos}, and quantified the effect of vertical displacement and oscillation for low-beam truck-to-truck links \cite{mohamed2023}. Adaptive field-of-view receivers further illustrate that directional receiver geometry is a practical design variable \cite{cailean2024}.

These references motivate directional geometry and measured optical-link variability, but they do not calibrate the analytical PC-FMCW-informed surrogate used below. The simulator latent vector remains a deliberately small model-mismatch family.

\section{Latent Link Model and Belief}
The modeled latent vector is
\begin{equation}
\phi=[\alpha_{\mathrm{loss}},\delta_{\mathrm{beam}},k_{\mathrm{angular}}],
\end{equation}
where $\alpha_{\mathrm{loss}}$ scales distance loss, $\delta_{\mathrm{beam}}$ is a modeled boresight offset, and $k_{\mathrm{angular}}$ scales angular attenuation. Simulator truth, the nominal vector, and the planner belief are distinct. C0--C3 never receive simulator truth.

C1--C3 maintain a discrete $3\times3\times3$ Bayesian grid. After an action is executed, the posterior is updated using the realized modeled-SNR observation and a Gaussian likelihood. The update is causal: the observation generated by an action cannot affect selection of that same action.

For candidate trajectory $\tau$, the predictive information proxy is
\begin{equation}
G(\tau)=\sum_j \frac{1}{2}\log\left(1+
\frac{\mathrm{Var}_{\phi}[\mu_z(\phi,\tau,j)]}{\sigma_z^2}\right).
\end{equation}
This is an approximate predictive-information score, not exact mutual information or a physically calibrated Fisher-information matrix.

\section{Decision-Triggered Planning}
Let $J(\tau,\phi)$ be candidate cost under latent hypothesis $\phi$, and let $\bar{\tau}$ minimize posterior-expected cost. Decision disagreement and expected decision regret are
\begin{align}
U_D &= P_{\phi}\!\left[\arg\min_{\tau}J(\tau,\phi)\neq\bar{\tau}\right],\\
R_D &= E_{\phi}\!\left[J(\bar{\tau},\phi)-\min_{\tau}J(\tau,\phi)\right].
\end{align}
C3 enables the information reward only when both declared gates are exceeded.

The planner family is:
\begin{itemize}
\item C0: nominal parameters, no update;
\item C1: passive Bayesian calibration, no information reward;
\item C2: unconditional information-seeking;
\item C3: decision-triggered information-seeking;
\item C4: true-parameter connectivity reference.
\end{itemize}

All planners share road, speed, static-obstacle, time-aligned dynamic-target, and stopping-feasibility filters. Information value is evaluated only after this common hard-safety filter and cannot make an infeasible candidate feasible.

\section{Frozen Protocol}
Development seeds were 31000--31019. The complete 27-setting grid was executed before deterministic selection. Setting \texttt{dev\_t010\_i025\_p020} was selected: decision threshold 0.10, information weight 0.25, probe weight 0.20, and minimum expected regret 0. The protocol was frozen at 2026-09-25 07:28:10 UTC before confirmatory seeds 32000--32049 were opened.

Six declared scenarios (A--F) are repeated within each seed. The independent inferential unit is the seed after averaging scenarios. Primary comparison families are C0--C1, C1--C2, C2--C3, and C3--C4. Deterministic paired bootstrap intervals, paired Wilcoxon tests, Cohen's $d_z$/rank-biserial effects, and Holm correction are reported across the seven declared metrics in each family.

\section{Confirmatory Results}
The confirmatory artifact contains 1,500 planner/scenario/seed episodes and 45,000 timestep rows. The sampled hard gate recorded zero collision episodes, zero no-candidate steps, and zero static-clearance-violation steps.

\begin{table}[t]
\caption{Frozen confirmatory planner means.}
\label{tab:means}
\centering
\small
\input{generated/table_planner_means.tex}
\end{table}

\begin{table}[t]
\caption{Primary seed-level cumulative-regret effects. Delta is planner B minus planner A.}
\label{tab:regret}
\centering
\small
\input{generated/table_primary_regret_effects.tex}
\end{table}

\begin{figure}[t]
\centering
\includegraphics[width=\columnwidth]{generated/fig_overall.pdf}
\caption{Frozen confirmatory planner means for cumulative decision regret (left) and probe fraction (right). C2 illustrates the cost of unconditional information seeking; C3 suppresses most of that probing burden.}
\label{fig:overall}
\end{figure}

Passive C1 lowers parameter error relative to C0. C2 learns parameters somewhat more aggressively than C1 but incurs much greater decision regret. C3 accepts higher parameter error than C2 while dramatically reducing regret and probing. Thus parameter learning and decision quality are not interchangeable endpoints.

C1 and C3 were not a predeclared primary pair. Descriptively, C1 has lower overall mean regret than C3, so no post-hoc C3-over-C1 superiority claim is made.

\section{Decision-Relevance Diagnostics}
Scenario E was designed so that parameter uncertainty should be decision irrelevant. C3 has zero mean probe fraction and zero mean regret, whereas C2 still probes (0.004667) and incurs regret (0.008610).

Scenario F was designed to be decision critical. Here C3 activates probing (0.003333) but does not improve regret over passive C1: mean regret is 0.031984 (C0), 0.028037 (C1), 0.140882 (C2), 0.034416 (C3), and 0 (C4). This negative result is retained rather than followed by post-confirmatory tuning.

\begin{figure}[t]
\centering
\includegraphics[width=\columnwidth]{generated/fig_decision_relevance.pdf}
\caption{Mechanism diagnostics for passive C1, unconditional-active C2, and decision-triggered C3. In scenario E, C3 suppresses probing entirely; in scenario F, C3 probes but does not outperform passive C1 on regret.}
\label{fig:decision-relevance}
\end{figure}

\section{Directional Mechanism Ablation}
A separate development-only ablation removes angular attenuation while preserving the planner/safety machinery and frozen-selected hyperparameters. Across 20 development seeds, mean regret is zero for C0/C4 and numerical roundoff for C1--C3; probe fraction is zero for every planner. C1 reduces the remaining distance-loss parameter error relative to C0, while angular latents are intentionally non-identifiable.

This collapse supports dependence on directional geometry inside the analytical model. It is not physical optical validation.

\section{Measured Support and Claim Boundaries}
The measured V-VLC support study uses the dataset associated with Turan and Coleri \cite{turan2021}. On held-out spatial groups, distance-only MAE is 5.587261 dB and directional MAE is 5.580011 dB. The directional-minus-distance-only absolute-error effect is -0.005031 dB (95\% CI [-0.019734,0.010134], Wilcoxon $p=0.375269$), effectively null evidence for a meaningful directional gain.

A separate CICV5G analysis uses the measured dataset of Zhang \emph{et al.} \cite{zhang2026}. It supports measured 5G QoS/pose prediction and route-constrained offline replay, not optical calibration or closed-loop real-vehicle validation.

\section{Discussion and Limitations}
The frozen result supports decision relevance as a guard against unnecessary information seeking. It does not establish that the implemented active probing policy is better than passive Bayesian calibration. Scenario F suggests that detecting decision relevance is insufficient by itself; limited probing actions, the short-horizon information proxy, or value scaling may matter, but the frozen study does not identify which explanation dominates.

The latent parameters and observation noise are modeled. The analytical link is not physically calibrated. C4 is a model-relative oracle, not evidence of real-world optimality. The candidate lattice constrains identifiability. The sampled hard gate is not a safety guarantee. The V-VLC directional comparison is null and CICV5G is a different radio modality.

\section{Conclusion}
Decision-triggered active self-calibration removes most of the large decision-regret and probing penalty produced by unconditional information seeking and suppresses probing when uncertainty is decision irrelevant. However, it does not demonstrate a confirmatory advantage over passive calibration. The supported conclusion is therefore deliberately narrow: decision relevance is useful for deciding when \emph{not} to probe, while the present probing policy does not yet show that active calibration improves downstream decisions beyond passive Bayesian updating.

\section*{Data and Code Availability}
Frozen protocol records, compact machine-readable summaries, artifact IDs/digests, workflows, and build instructions are included in the repository. Full Actions artifacts remain the authoritative raw evidence.

\section*{Declarations}
Author affiliation/postal address, institutional email, short biography, ORCID, acknowledgements/funding, and conflict-of-interest disclosure must be confirmed before any venue portal submission. T-IV submission additionally requires resolution of its public-repository checklist item for the already-public manuscript history.

\bibliographystyle{IEEEtran}
\bibliography{references}
\end{document}
